\section{2020-2021学年线性代数II（H）小测答案}

\begin{center}
    任课老师：刘康生\hspace{4em} 考试时长：45分钟
\end{center}
\begin{enumerate}
    \item 复内积空间上，以下条件等价：
    \[
        T\text{ 正规}\iff T^*T=TT^*
        \iff \|Tv\|=\|T^*v\|\quad(\forall v)
    \]
    \[
        \iff V\text{ 有一组由 }T\text{ 的特征向量组成的规范正交基}
        \iff T\text{ 可酉对角化}.
    \]
    还可等价表述为：$T$ 的每个不变子空间的正交补仍不变；不同特征值的特征子空间相互正交；$T$ 与 $T^*$ 可同时酉对角化.

    \item 实内积空间上，$T$ 正规等价于 $T^*T=TT^*$，也等价于 $\|Tv\|=\|T^*v\|$ 对所有 $v$ 成立. 还等价于存在一组规范正交基，使 $T$ 的矩阵是若干个一维实块 $[a]$ 与二维块
    \[
        \begin{pmatrix}a&-b\\ b&a\end{pmatrix}\quad(b\neq 0)
    \]
    的正交直和. 对实矩阵 $A$，还等价于 $A^{\mathrm{T}}A=AA^{\mathrm{T}}$，或等价于其对称部分与反对称部分可交换.

    \item 在有限维情形下存在. 由 $\|Tv\|=\|Sv\|$ 知 $\ker S=\ker T$. 定义
    \[
        U_0:\im S\to \im T,\qquad U_0(Sv)=Tv.
    \]
    若 $Sv_1=Sv_2$，则 $v_1-v_2\in\ker S=\ker T$，故定义良好. 且
    \[
        \|U_0(Sv)\|=\|Tv\|=\|Sv\|,
    \]
    所以 $U_0$ 是 $\im S$ 到 $\im T$ 的等距同构. 又 $\dim(\im S)^\perp=\dim(\im T)^\perp$，将 $U_0$ 任意延拓为整个 $V$ 上的等距同构 $U$，即有 $T=US$.

    \item
    \begin{enumerate}
        \item 直接计算得
        \[
            A^{\mathrm{T}}A=AA^{\mathrm{T}}=2I,
        \]
        故 $A$ 正规.

        \item 取
        \[
            p_1=\left(\frac12,\frac1{\sqrt2},\frac12,0\right)^{\mathrm{T}},\quad
            p_2=\left(\frac12,-\frac1{\sqrt2},\frac12,0\right)^{\mathrm{T}},
        \]
        \[
            p_3=(0,0,0,1)^{\mathrm{T}},\quad
            p_4=\left(\frac1{\sqrt2},0,-\frac1{\sqrt2},0\right)^{\mathrm{T}}.
        \]
        它们构成规范正交基，且
        \[
            Ap_1=\sqrt2p_1,\quad Ap_2=-\sqrt2p_2,\quad
            Ap_3=\sqrt2p_4,\quad Ap_4=-\sqrt2p_3.
        \]
        令 $P=(p_1,p_2,p_3,p_4)$，则
        \[
            P^{\mathrm{T}}AP=
            \begin{pmatrix}
                \sqrt2&0&0&0\\
                0&-\sqrt2&0&0\\
                0&0&0&-\sqrt2\\
                0&0&\sqrt2&0
            \end{pmatrix}.
        \]
    \end{enumerate}
\end{enumerate}

\clearpage
